% !TEX root = ../bb_AiRa_Kappaleet.tex

%% HARRIS: Chapter 10

% ö = \%c3\%b6
% ä = \%C3\%A4

\xdef\sectiontitle{\IIsectiontitle} %aiheuttama
\TOCrefs

%%%%%%%%%%%%%%%
\begin{frame}[label=2_10_a]%[label=2_9_TOC1]
\frametitle{\texorpdfstring{\culine[\sectiontitleulinecolor]{\sectiontitle}}{\sectiontitle} }
\label{\TOCname}

\vspace{\chapterTOCvksip}
\begin{itemize}[leftmargin=0.5cm,itemsep=3mm]
    \item[2.1]<1-> \hyperlink{2_1_a}{\IIaaSubsectiontitle}
    \item[2.2]<1-> \hyperlink{2_2_a}{\IIaSubsectiontitle}
    \item[2.3]<1-> \hyperlink{2_3_a}{\IIbSubsectiontitle}	
    \item[2.4]<1-> \hyperlink{2_4_a}{\IIcSubsectiontitle}	
    \item[2.5]<1-> \hyperlink{2_5_a}{\IIdSubsectiontitle}
    \item[2.6]<1-> \hyperlink{2_6_a}{\IIeSubsectiontitle}
    \item[2.7]<1-> \hyperlink{2_7_a}{\IIfSubsectiontitle}
    \item[2.8]<1-> \hyperlink{2_8_a}{\IIgSubsectiontitle}
    \item[2.9]<1-> \hyperlink{2_9_a}{\IIhSubsectiontitle}
    \item[2.10]<1-> \hyperlink{2_10_a}{\CDAlert[blue]{\IIiSubsectiontitle}}
\end{itemize}

%\endofslide 
\end{frame}
%%%%%%%%%%%%%%%

\xdef\sectiontitle{\IIiSubsectiontitle} 

%%%%%%%%%%%%%%%
\begin{frame}
\frametitle{\texorpdfstring{\culine[\sectiontitleulinecolor]{\sectiontitle}}{\sectiontitle} }
\framesubtitle{Hiilen allotrooppien sovelluskohteita}
\appear{}
%On fullerenes and Carbon nanotubes (CNT)
\begin{itemize}
	\item Materiaalitieteen kehitys on johtanut lukuisiin hiilestä koostuviin \appear{3D,} 
\end{itemize}
%\vspace{1.6mm}
\begin{minipage}{6.75cm}
\begin{itemize}[itemsep=1.6mm]
	\item[] 2D\appear{ ja jopa 1D} \appear{nanorakenteisiin} \appear{kuten fullereeneihin}\appear{, \CDAlert[red]{grafeeniin}}\appear{, ja \CDAlert[blue]{hiilinanoputkiin} }
	
	\item Hiilinanoputkien \appear{ominaisuuksia voidaan parantaa seos\-ta\-mal\-la niitä funktio\-naa\-li\-sil\-la atomeilla} \appear{(\CDAlert[blue]{optoelektroniikka})}

\item Fullereenejä on käytetty myös molekylaarisina kuulalaakereina\appear{, mik\-roskooppisina siivilöinä} \appear{ja mit\-ta\-tik\-kuina}

\item Hiilinanoputkilla ja grafeenilla on erittäin suuri \Alert[fuchsia]{mekaaninen kestävyys} 
\implies{ Urheiluvälineet ja suojavaatteet }
	\end{itemize}
\end{minipage}

\xdef\loc{south east}
\begin{tikzpicture}[remember picture,overlay] % anchor=south
    \node (A) [yshift=-0.25*\figYshift,xshift=\figXshift, anchor=\loc] at (current page.\loc) {\includegraphics[page=1,width=7cm]{Figures_booklet/SOM_10_Bonding_figures/SOM_Bonding_Figure_57.png}}; %scale=\figscale. % 8cm max height
\end{tikzpicture}
\footnote{\HarrisCR[1-]}

%% 58 and 61 figs too
\endofslide 
%\endofsection
\end{frame}
%%%%%%%%%%%%%%%

%%%%%%%%%%%%%%%
\begin{frame}
\frametitle{\texorpdfstring{\culine[\sectiontitleulinecolor]{\sectiontitle}}{\sectiontitle} }
\framesubtitle{Hiilinanoputket (carbon nanotubes, CNT)}
\appear{}
\begin{itemize}
	\item Erikoinen ryhmä akiraalisia yksiseinäisiä hiilinanoputkia on \Alert[red]{metallisia}\appear{, mutta loput ovat joko \Alert[blue]{pienen tai keskisuuren energiaraon puolijohteita}}\appear{, jolloin niillä on potentiaalia nanoelektroniikassa} \appear{(\CDAlert[fuchsia]{kanavatransistorit} ja muut komponentit)} \appear{ja esim. valolle herkät nanoantennit (analogia radioantenneihin)}
\end{itemize}

\xdef\loc{south west}
\begin{tikzpicture}[remember picture,overlay] % anchor=south
    \node (A) [yshift=-0.25*\figYshift,xshift=-\figXshift, anchor=\loc] at (current page.\loc) {\includegraphics[page=1,width=7cm]{Figures_booklet/SOM_10_Bonding_figures/SOM_Bonding_Figure_59.png}}; %scale=\figscale. % 8cm max height
\end{tikzpicture}

\xdef\loc{south east}
\begin{tikzpicture}[remember picture,overlay] % anchor=south
    \node (A) [yshift=-0.25*\figYshift,xshift=\figXshift, anchor=\loc] at (current page.\loc) {\includegraphics[page=1,width=7cm]{Figures_booklet/SOM_10_Bonding_figures/SOM_Bonding_Figure_60.png}}; %scale=\figscale. % 8cm max height
\end{tikzpicture}
\footnote{\HarrisCRshort[1-]}

\endofslide 
%\endofsection
\end{frame}
%%%%%%%%%%%%%%%



